\documentclass{article}
\usepackage{amsmath,amssymb,amsthm,hyperref}
\newtheorem{theorem}{Theorem}
\begin{document}
\section*{Complete classification of unrounded packing equality through order four}
Let $\mathcal C\subset\operatorname{Sym}^N(\mathbb C^2)$ be a nonzero
$t$-anticoherent subspace of dimension $k$, with $N\ge t\ge1$.
The unrounded rectangular-tensor packing inequalities at orders one
through four are
\[
2k\le N,\qquad4k\le N+1,\qquad6k\le N,\qquad9k\le N+1,
\]
respectively. Equality here refers to these unrounded inequalities,
not to their integer-rounded versions or to the strongest available
upper bound.

\begin{theorem}
Equality in these four inequalities is possible precisely as follows:
\[
\begin{array}{c|c|c}
 t&\text{equality}&\text{all possible }(N,k)\\\hline
 1&2k=N&(2,1),(4,2),(6,3)\\
 2&4k=N+1&(7,2)\\
 3&6k=N&(6,1)\\
 4&9k=N+1&\text{none}.
\end{array}
\]
Consequently the following strict refinements hold:
\[
\boxed{6K(N,3)\le N-1\quad(N\ge7),}
\qquad
\boxed{9K(N,4)\le N\quad(N\ge4).}
\]
These do not assert that the refined or integer-rounded bounds are attained.
\end{theorem}

\paragraph{Order one.}
The already proved auxiliary-spin bound $3k\le N+3$ implies $N\le6$
when $2k=N$. Since $N$ is a positive even integer, only $N=2,4,6$
remain. Explicit first-order bases, with normalized Dicke states, are
\begin{align*}
N=2:&\quad (|D_2^{(0)}\rangle+|D_2^{(2)}\rangle)/\sqrt2,\\
N=4:&\quad (|D_4^{(0)}\rangle+|D_4^{(4)}\rangle)/\sqrt2,
                   \quad |D_4^{(2)}\rangle,\\
N=6:&\quad \sqrt{3/5}|D_6^{(1)}\rangle+\sqrt{2/5}|D_6^{(6)}\rangle,\\
&\quad \sqrt{2/5}|D_6^{(0)}\rangle-\sqrt{3/5}|D_6^{(5)}\rangle,
                   \quad |D_6^{(3)}\rangle.
\end{align*}
Their compressed spin generators vanish identically.

\paragraph{Order two.}
The bound $5k\le N+5$ implies $N\le15$ at equality
$4k=N+1$. The candidates are $N=3,7,11,15$.
A pure second-order state needs $N\ge4$, excluding $N=3$.
For every $N\ge8$ the already established auxiliary-spin-$3/2$ theorem
$4k\le N-2$ excludes equality, hence excludes $N=11,15$.
At $N=7$ an exact optimal plane has basis
\[
\sqrt{3/10}|D_7^{(0)}\rangle+\sqrt{7/10}|D_7^{(5)}\rangle,
\qquad
\sqrt{7/10}|D_7^{(2)}\rangle-\sqrt{3/10}|D_7^{(7)}\rangle.
\]
For these vectors all degree-two compressed spin moments agree with the
normalized physical trace. The auxiliary-spin theorems used here are in
\href{../t2-spin2-inertia/proof.tex}{the spin-two proof} and
\href{../ancilla-inertia/proof.tex}{the spin-three-halves proof}.

\paragraph{Order three.}
The new bound $7k\le N+7$ implies $N\le42$ when $6k=N$.
Thus only $N=6,12,18,24,30,36,42$ remain.
For $N\ge7$, equality in the rectangular-tensor channel requires a
nonscalar zero multiplier below the highest physical tensor rank.
Here the highest rank $L=N$ is automatically annihilated because the
rectangular channel outputs spin $(N-1)/2$. If it were the only
nonscalar kernel rank, the projector would have the form
$P=\alpha I+X_N$. The
\href{../general-bounds/proof.tex}{highest-rank projector obstruction}
then vectorizes $X_N$ into a pure symmetric $2N$-qubit state of AC
order $N-1$. Its pure-state packing inequality would require
$(N+1)^2\le8N+1$, hence $N\le6$, a contradiction.
Thus, for the remaining $N\ge12$ candidates, the necessary zero is
equivalent to
\[
2L(L+1)=(N-1)(N+2),\qquad 4\le L\le N-1.
\]
An exact check of the finite list leaves only $(N,L)=(6,4)$.
At $N=6$ the octahedral state
\[
(|D_6^{(1)}\rangle+|D_6^{(5)}\rangle)/\sqrt2
\]
is third-order AC. The auxiliary-spin theorem therefore turns the
previous Pell-type necessary condition into a complete equality
classification: all larger Pell solutions are outside the possible
range $N\le42$. In particular even the next arithmetic candidate,
$N=78$, cannot attain $6k=N$.

\paragraph{Order four.}
The new bound $10k\le N+10$, valid for $N\ge9$, gives $N\le80$
when $9k=N+1$. There is also the smaller candidate $N=8$.
Hence the complete candidate list is
\[
N=8,17,26,35,44,53,62,71,80.
\]
Under packing equality, the normalized tensor channel
\[
\mathcal E_2(A)=\sum_{\ell=0}^2\sum_{m=-\ell}^{\ell}
 T_{\ell m}A T_{\ell m}^{\dagger}
\]
satisfies $\mathcal E_2(P)=I/(N+1)$ for the code projector $P$.
Its multiplier on rank $L$ is
\begin{align*}
e_{2,L}
&=\sum_{\ell=0}^2(-1)^{N+\ell+L}(2\ell+1)
 \begin{Bmatrix}N/2&N/2&\ell\\N/2&N/2&L\end{Bmatrix}\\
&=\frac{3[10x^2+(16-12Q)x+3Q(Q-3)]}{(N+1)Q(Q-3)},
\qquad x=L(L+1),\quad Q=N(N+2).
\end{align*}
An exact independent evaluation of every admissible $5\le L\le N$
for this finite candidate list gives no zero, except $L=16$ at $N=17$.
Absence of a zero would force $P=kI/(N+1)$, contradicting its being a
nontrivial projector.

The remaining $N=17$ case is excluded by the existing
\href{../n17-t4/proof.tex}{rank-sixteen vectorization proof}.
For completeness, the key chain is as follows. Equality forces
$P=I_{18}/9+X_{16}$. First-order AC gives a rank-four projector
$Q'=2\operatorname{Tr}_1P=4I_{17}/17+Y_{16}$.
Then
\[
Y_{16}^2=\frac9{17}Y_{16}+\frac{52}{289}I_{17},
\qquad \operatorname{Tr}Y_{16}^2=\frac{52}{17}.
\]
Time-reversed vectorization of $Y_{16}$ produces a pure symmetric
32-qubit state. Its sixteen-qubit reduction equals
$I_{17}/17+9Y_{16}/52$; tracing one more qubit annihilates rank sixteen.
The resulting pure state would therefore be fifteenth-order AC.
Its tensor-Gram bound requires $64\le33$, a contradiction.
Thus the last candidate also fails.

\paragraph{Independent exact verification.}
The companion script imports no catalogue bound implementation.
It reconstructs Wigner-$6j$ channel multipliers directly, checks their
zeros and the fourth-order rational formula for every candidate, verifies
the finite first- and second-order eliminations, and checks all ordered
spin-word compressions for the displayed attaining witnesses exactly.
The unbounded ranges are excluded analytically by the auxiliary-spin
inequalities; no finite numerical scan is used to infer them.
\end{document}
